\documentclass[11pt]{article}
% ---------------------------------------------------------------------------
% Results are read from the experiment harness (cvebench). Point \ResultsDir at
% a run directory; tables/figures that do not exist yet render as red placeholders.
%   ../results/main  -> real run      ../results/mock -> pipeline dry run (watermarked)
\newcommand{\ResultsDir}{../results/main}
% ---------------------------------------------------------------------------
\usepackage[a4paper,margin=2.6cm]{geometry}
\usepackage[T1]{fontenc}
\usepackage{lmodern,microtype}
\usepackage{amsmath,amssymb,amsthm}
\usepackage{graphicx,booktabs,multirow,adjustbox}
\usepackage{xcolor,etoolbox,enumitem}
\usepackage[round,authoryear]{natbib}
\usepackage{algorithm,algpseudocode}
\usepackage[most]{tcolorbox}
\usepackage[hidelinks]{hyperref}
\usepackage[mathlines]{lineno}
\linenumbers

\IfFileExists{\ResultsDir/tables/provenance.tex}{\input{\ResultsDir/tables/provenance.tex}}{\newcommand{\ResultsProvenance}{NONE}}
\ifdefstring{\ResultsProvenance}{MOCK}{%
  \usepackage[text={MOCK DATA -- NOT RESULTS},scale=0.45,color=red!22]{draftwatermark}}{}

\newtheorem{definition}{Definition}[section]
\newtheorem{proposition}{Proposition}[section]
\newtheorem{remark}{Remark}[section]
\newtcolorbox{keyresult}{colback=black!4,colframe=black!55,boxrule=0.6pt,arc=1pt,
  title=Key Result,fonttitle=\bfseries\small,fontupper=\small}
\newcommand{\TBD}[1]{\textcolor{red}{[\textbf{TBD:} #1]}}
\newcommand{\Fab}{\mathrm{Fab}}\newcommand{\Acc}{\mathrm{Acc}}\newcommand{\Abs}{\mathrm{Abs}}
\newcommand{\Cov}{\mathrm{Cov}}
\newcommand{\resultstable}[3]{%
  \begin{table}[htbp]\centering\caption{#2}\label{#3}\small
  \IfFileExists{\ResultsDir/tables/#1.tex}%
    {\begin{adjustbox}{max width=\textwidth}\input{\ResultsDir/tables/#1.tex}\end{adjustbox}}%
    {\fbox{\parbox{0.92\textwidth}{\centering\vspace{1.2em}\textcolor{red}{Table \texttt{\detokenize{#1}.tex} is generated by \texttt{python -m cvebench.cli analyze}}\vspace{1.2em}}}}%
  \end{table}}
\newcommand{\resultsfig}[4][0.95]{%
  \begin{figure}[htbp]\centering
  \IfFileExists{\ResultsDir/figures/#2.pdf}%
    {\includegraphics[width=#1\textwidth]{\ResultsDir/figures/#2.pdf}}%
    {\fbox{\parbox{0.92\textwidth}{\centering\vspace{3em}\textcolor{red}{Figure \texttt{\detokenize{#2}.pdf} is generated by \texttt{python -m cvebench.cli analyze}}\vspace{3em}}}}%
  \caption{#3}\label{#4}\end{figure}}

\title{Does Consensus Prevent Confabulation? Parallel Multi-LLM Ensembles versus Single Models for Hallucination-Safe CVE Record Analysis}
\author{Blinded manuscript}
\date{}

\begin{document}
\maketitle
\ifdefstring{\ResultsProvenance}{MOCK}{\begin{center}\fcolorbox{red}{red!8}{\parbox{0.9\textwidth}{\centering\small\textbf{Pipeline dry run.} All tables and figures in this build come from simulated models on synthetic records and must not be read as findings.}}\end{center}}{}

\begin{abstract}
Large language models (LLMs) are now used to triage Common Vulnerabilities and Exposures (CVE) records: they assign CVSS scores, map weaknesses to CWE categories and identify affected products, and their output feeds patch-prioritisation and compliance pipelines in which a single fabricated value propagates silently. Prior work has shown that retrieval grounding and post-generation verification can eliminate fabrication when an authoritative copy of the answer exists. Many operational cases have no such copy: newly published CVEs awaiting enrichment, or fields that must be inferred from the advisory text. A natural engineering response is to query several heterogeneous LLMs in parallel and trust only what they agree on. Whether this works depends on a question that is rarely measured: do different LLMs fabricate \emph{the same} wrong values?

We present a controlled, paired comparison of parallel multi-LLM ensembles against their constituent single models on real National Vulnerability Database (NVD) records. Two tasks are evaluated on identical queries: \emph{parametric recall} (CVE identifier only; CVSS v3.1 base score and CWE), on a post-training-cutoff stratum and a pre-cutoff control stratum, and \emph{record-grounded inference} (identifier and description; CVSS vector and score, CWE, vendor, product). Every field is graded as correct, fabricated or abstained, extending the umbrella hallucination rate $\mathcal{H}$ of our earlier study. We add three hallucination flags that need no ground truth: identifiers that do not exist, names absent from the input record, and scores inconsistent with the model's own CVSS vector. We compare plurality, strict-majority, unanimous, \emph{verified-majority}, metric-wise CVSS voting and LLM-judge fusion against the best single model, selected on a held-out split. We also compare against a compute-matched self-consistency baseline that spends the same number of calls on one model, which separates the effect of model \emph{diversity} from that of extra sampling. We prove that strict-majority voting suppresses a fabricated value at a rate exponential in ensemble size only if member models rarely produce \emph{that specific} wrong value. We also prove that verification before voting can only trade coverage for safety, never the reverse. The bound motivates a new diagnostic, the identical-fabrication rate. \TBD{headline results: fabrication of best single vs.\ best ensemble on T1 and T2 with 95\% CIs; accuracy cost; identical-fabrication rate; McNemar $p$; cost multiplier}. The results give security teams a quantitative basis for deciding when a multi-model ensemble is worth its cost, and which aggregation rule to deploy.
\end{abstract}

\noindent\textbf{Keywords:} large language models; hallucination; ensemble methods; CVE; CVSS; CWE; vulnerability management; selective prediction

% ============================================================================
\section{Introduction}\label{sec:intro}

\paragraph{The fabrication problem in vulnerability pipelines.}
LLM components are increasingly placed at the front of vulnerability-management pipelines, where they read advisories and emit structured metadata (a CVSS v3.1 base score, a CWE category, an affected vendor and product) that downstream systems treat as data \citep{fayyazi2025proverag,alam2024ctibench}. A fabricated value in this setting is not a stylistic defect. It is a syntactically valid record that passes schema validation and changes patch ordering, service-level timers and risk scores. Our earlier study \citep{anon2026grounding} measured this failure mode directly on $N{=}1000$ NVD records. Unassisted models fabricated or refused on at least 94.2\% of score-recall queries. Retrieval of the authoritative NVD record, followed by a dual-gate exact-match and entailment check, drove fabrication to zero for two open-weight models.

\paragraph{When there is nothing to retrieve.}
Retrieval grounding presupposes that the answer already exists in an authoritative store. This is often not the case in practice. Newly published CVEs can wait for NVD enrichment, and in that interval no authoritative score, CWE or product configuration exists to retrieve. CNA-supplied descriptions are often the only input, and analysts must \emph{infer} metadata from them \citep{shahid2021cvssbert,jafarikhah2025description}. Here, retrieval-then-verify reduces to verifying a guess against nothing. The question becomes how to make inference itself less likely to fabricate.

\paragraph{Parallel ensembles as a candidate safeguard.}
A widely used engineering answer is redundancy: send the same record to several independently trained LLMs in parallel, and release a value only when enough of them agree. The intuition is the classical ensemble argument \citep{dietterich2000}. If members err independently, a wrong value rarely wins a vote. Recent work applies this idea to vulnerability \emph{detection} in source code, reporting 10--12\% higher accuracy than the average single model \citep{zibaeirad2025diverse}. In contrast, an ensemble meta-classifier over CVSS predictions from descriptions gave only marginal gains, because the models misclassified many of the same CVEs \citep{jafarikhah2025description}. A large-scale study of more than 350 LLMs finds that models agree on the same wrong answer far more often than chance, and that larger, more accurate models are more correlated \citep{kim2025correlated}. If CVE fabrications are similarly shared (for example, the same ``typical'' score of 9.8 or the same generic CWE-20), voting will ratify them instead of filtering them.

\begin{keyresult}
\textbf{The core question.} Given the same CVE record, does a parallel ensemble of heterogeneous LLMs return fewer fabricated values than the best single LLM, and at what cost in coverage, compute and latency? Is any improvement due to model \emph{diversity}, or merely to spending more calls?
\end{keyresult}

No prior study answers this for CVE records with fabrication-specific metrics. HalluCVE measures hallucination in LLM-generated CVE content but evaluates single models only \citep{dao2026hallucve}. DVDR-LLM ensembles models for code-level detection and repair but does not measure fabrication \citep{zibaeirad2025diverse}. CVSS-from-description studies report accuracy without separating fabrication from abstention \citep{jafarikhah2025description}.

\subsection{Research Questions}\label{sec:rq}
\begin{itemize}[leftmargin=2.2em]
\item[\textbf{RQ1.}] \emph{Parametric recall.} Given only a CVE identifier, does a parallel ensemble fabricate fewer CVSS scores and CWE identifiers than the best single model, on records published after the models' training cutoffs (where the only safe answer is abstention) and on a pre-cutoff control?
\item[\textbf{RQ2.}] \emph{Record-grounded inference.} Given the CVE description, does an ensemble improve accuracy and reduce fabrication for the CVSS vector, score, CWE, vendor and product, and does it reduce ungrounded and self-inconsistent outputs?
\item[\textbf{RQ3.}] \emph{Diversity versus compute.} Is the ensemble effect explained by model diversity, or is it matched by self-consistency over one model with the same number of calls? How often do different models fabricate \emph{identical} values?
\item[\textbf{RQ4.}] \emph{Aggregation and operating point.} Which aggregation rule (plurality, strict majority, unanimity, verified majority, metric-wise voting, LLM-judge fusion) gives the best fabrication--coverage trade-off, how does it scale with ensemble size, and at what cost?
\end{itemize}

\paragraph{Hypotheses.} We pre-register five hypotheses. \textbf{H1}: strict-majority ensembles have lower fabrication than the best single model on both tasks, at a cost in coverage. \textbf{H2}: plurality voting does \emph{not} reduce fabrication on post-cutoff recall, because members converge on a few popular wrong values. \textbf{H3}: at matched call budgets, heterogeneous ensembles have lower identical-fabrication rates, and therefore lower post-vote fabrication, than self-consistency. \textbf{H4}: deterministic verification before voting reduces fabrication further, without any increase in fabrication (guaranteed by Proposition~\ref{prop:verified}). \textbf{H5}: LLM-judge fusion introduces values that no member proposed at a non-zero rate.

\subsection{Contributions}
\begin{enumerate}[leftmargin=2.2em]
\item A paired evaluation design that compares parallel ensembles, their constituent single models and a compute-matched self-consistency baseline on identical NVD queries. The best single model is selected on a disjoint development split, which avoids the optimistic ``best-of-$n$'' reference.
\item A field-level outcome taxonomy (correct / fabricated / abstained) with three ground-truth-free hallucination flags (\emph{invalid}, \emph{ungrounded}, \emph{inconsistent}) for CVE metadata, including an exact CVSS v3.1 recomputation to detect self-contradictory vector/score pairs.
\item Two propositions: a Hoeffding bound showing that majority voting suppresses a fabricated value only to the extent that members rarely emit \emph{that specific} value, and a monotonicity result showing that pre-vote verification can only reduce fabrication. Together they motivate the \emph{identical-fabrication rate} as the diagnostic that predicts whether an ensemble will help.
\item An open, resumable harness (\texttt{cvebench}) supporting OpenAI-compatible, Anthropic, local Hugging Face and vLLM backends. It logs every raw response and generates every table and figure in this paper from those logs. \TBD{repository URL at camera-ready}.
\end{enumerate}

\subsection{Paper Organisation}
Section~\ref{sec:related} reviews related work. Section~\ref{sec:math} formalises outcomes, aggregation operators and the two propositions. Section~\ref{sec:method} describes the corpus, model registry and pipeline algorithms. Section~\ref{sec:results} reports results, Section~\ref{sec:discussion} discusses them, Sections~\ref{sec:limitations}--\ref{sec:future} give limitations and future work, and Section~\ref{sec:conclusion} concludes.

% ============================================================================
\section{Related Work}\label{sec:related}

\subsection{Hallucination in LLMs}
Hallucination, meaning fluent output unsupported by the input or by world knowledge, is surveyed by \citet{ji2023survey} and \citet{huang2025survey}. \citet{kalai2025hallucinate} argue that standard training and evaluation reward guessing over abstaining. This argument applies directly to structured security metadata, where a guessed score is always ``well formed''. Black-box detectors exploit inconsistency across samples \citep{manakul2023selfcheckgpt} or semantic entropy over meaning clusters \citep{farquhar2024semantic}. \citet{kadavath2022know} show that models can partly predict their own correctness. Our setting differs in that ground truth is exact and machine-checkable, so fabrication can be measured rather than estimated.

\subsection{LLMs for CVE Metadata}
CVSS prediction from descriptions predates LLMs \citep{shahid2021cvssbert}. CTIBench evaluates root-cause (CVE-to-CWE) mapping and severity prediction \citep{alam2024ctibench}. A fine-tuned 125M-parameter RoBERTa model is competitive with LLMs on CVE-to-CWE classification \citep{mosievskiy2026roberta}. \citet{jafarikhah2025description} find that LLMs misclassify many of the same CVEs and that ensemble meta-classifiers help only marginally. \citet{ullah2024llms} show that LLMs cannot yet reliably identify and reason about vulnerabilities in code. HalluCVE reports pervasive hallucination in LLM-generated CVE content, including on post-cutoff vulnerabilities \citep{dao2026hallucve}. ProveRAG adds provenance and self-critique to retrieval-augmented vulnerability analysis \citep{fayyazi2025proverag}. Our earlier study isolates score recall and shows that retrieval plus dual-gate verification eliminates fabrication when the record is retrievable \citep{anon2026grounding}.

\subsection{Ensembles, Self-Consistency and Multi-Agent LLMs}
Classical ensemble theory ties ensemble gains to member diversity \citep{dietterich2000,kuncheva2003}. For LLMs, self-consistency samples one model several times and takes the majority answer \citep{wang2023selfconsistency}. Multi-agent debate \citep{du2024debate}, LLM-Blender \citep{jiang2023blender} and Mixture-of-Agents \citep{wang2024moa} combine several models. Cascades such as FrugalGPT trade accuracy against cost \citep{chen2023frugalgpt}. LLM-as-a-judge evaluation \citep{zheng2023judge} underlies judge-based fusion. \citet{kim2025correlated} measure correlated errors across hundreds of LLMs and warn against assuming diversity. In security, DVDR-LLM applies consensus thresholds to code-level vulnerability detection \citep{zibaeirad2025diverse}.

\subsection{Selective Prediction}
Abstention-aware evaluation is formalised in selective classification \citep{elyaniv2010selective,geifman2017selective}, where systems are compared by risk at a given coverage. Agreement thresholds in voting ensembles are a natural selection function. We therefore report risk--coverage curves rather than a single accuracy number.

\paragraph{Gap.} Prior CVE studies evaluate single models, and prior ensemble studies evaluate code-level detection or report accuracy only. None measures, on CVE records, whether parallel heterogeneous ensembles reduce \emph{fabrication}. None separates diversity from compute with a matched self-consistency baseline, or measures how often different models fabricate identical values.

% ============================================================================
\section{Mathematical Framework}\label{sec:math}

Table~\ref{tab:notation} summarises notation.

\begin{table}[htbp]\centering\small
\caption{Notation.}\label{tab:notation}
\begin{adjustbox}{max width=\textwidth}\begin{tabular}{@{}ll@{}}
\toprule
Symbol & Meaning \\ \midrule
$\mathcal{D}=\{R_i\}_{i=1}^{N}$ & corpus of CVE records $R_i=(c_i,d_i,\{y^*_{i,f}\})$: identifier, description, gold labels \\
$\mathcal{F}_{T}$ & graded fields of task $T$: $\mathcal{F}_{\mathrm{T1}}=\{\text{score},\text{cwe}\}$, $\mathcal{F}_{\mathrm{T2}}=\{\text{vec},\text{score},\text{cwe},\text{vendor},\text{product}\}$ \\
$\mathcal{M}=\{M_1,\dots,M_n\}$ & ensemble members (heterogeneous LLMs) \\
$\hat y_{m,i,f}$ & normalised answer of $M_m$ for field $f$ of record $i$; $\bot$ denotes abstention \\
$\mathcal{V}_f,\ \mathcal{X}_f$ & valid value space of field $f$; malformed or non-existent values \\
$o(\hat y,y^*)$ & outcome in $\{\mathsf{C},\mathsf{F},\mathsf{A}\}$ (correct, fabricated, abstained), Eq.~\eqref{eq:outcome} \\
$\mathsf{Agg}$ & aggregation operator, $\mathsf{Agg}:(\mathcal{V}_f\cup\mathcal{X}_f\cup\{\bot\})^n\to\mathcal{V}_f\cup\mathcal{X}_f\cup\{\bot\}$ \\
$\kappa_f(v)$ & vote count of value $v$, Eq.~\eqref{eq:count} \\
$\delta$ & numeric tolerance for scores, $0.01$ (as in \citealp{anon2026grounding}) \\
$S,\ \tau$ & self-consistency samples per model ($S{=}n$) and sampling temperature ($0.7$) \\
$\mathrm{CVSS}(\cdot)$ & CVSS v3.1 base-score function of a vector \citep{first2019cvss} \\
\bottomrule
\end{tabular}\end{adjustbox}
\end{table}

\subsection{Normalisation and Outcomes}
Model replies are parsed as JSON objects, with a regex fallback identical to the extractor of \citet{anon2026grounding} for score-only replies. Each field is mapped to a canonical value. Scores are rounded to one decimal and must lie in $[0,10]$. Vectors are canonicalised to \texttt{CVSS:3.1/AV:$\cdot$/\dots/A:$\cdot$} and must contain all eight base metrics with legal values. CWE identifiers are normalised to \texttt{CWE-}$k$ and checked against the MITRE catalogue \citep{mitre2024cwe}. Vendor and product names are lower-cased slugs. Explicit \texttt{null} and refusal tokens map to $\bot$. Values that are present but malformed or non-existent are kept as members of $\mathcal{X}_f$ so that they count as fabrications, not abstentions.

\begin{definition}[Field outcome]\label{def:outcome}
For a prediction $\hat y$ and gold label $y^*$ of field $f$,
\begin{equation}\label{eq:outcome}
o(\hat y,y^*)=\begin{cases}
\mathsf{A} & \hat y=\bot,\\
\mathsf{C} & \hat y\in\mathcal{V}_f\ \wedge\ \mathrm{match}_f(\hat y,y^*),\\
\mathsf{F} & \text{otherwise,}
\end{cases}
\end{equation}
where $\mathrm{match}_f$ is $|\hat y-y^*|\le\delta$ for scores, exact equality of all eight metrics for vectors, membership in the NVD-assigned CWE set for CWE (records labelled only \texttt{NVD-CWE-noinfo}/\texttt{Other} are not graded), and normalised equality or containment (length $\ge 4$) against any CPE vendor/product for names (records without CPE configurations are not graded).
\end{definition}

Over the $N_f$ gradable records of a field, a system's rates are
\begin{equation}\label{eq:rates}
\Acc=\tfrac{1}{N_f}\textstyle\sum_i \mathbf{1}[o_i{=}\mathsf{C}],\quad
\Fab=\tfrac{1}{N_f}\sum_i \mathbf{1}[o_i{=}\mathsf{F}],\quad
\Abs=\tfrac{1}{N_f}\sum_i \mathbf{1}[o_i{=}\mathsf{A}],\quad
\Cov=1-\Abs,
\end{equation}
with $\Acc+\Fab+\Abs=1$. The umbrella hallucination rate of \citet{anon2026grounding} is $\mathcal{H}=\Fab+\Abs=1-\Acc$. Following that study's Remark~3.1, we never report $\mathcal{H}$ without its decomposition, because an abstention is a safe, auditable non-answer whereas a fabrication silently corrupts downstream records. For selective-prediction comparisons we use selective accuracy $\Acc/\Cov$ and selective risk $\Fab/\Cov$ \citep{geifman2017selective}. Multi-field results are macro-averaged over fields.

\subsection{Ground-Truth-Free Hallucination Flags}
Three predicates detect hallucination without consulting NVD labels, so they can be deployed at inference time:
\begin{align}
\mathrm{Inv}(\hat y_f) &= \mathbf{1}[\hat y_f\in\mathcal{X}_f], \label{eq:inv}\\
\mathrm{Ungr}(\hat y_f,d) &= \mathbf{1}[f\in\{\text{vendor},\text{product}\}\wedge \hat y_f\neq\bot\wedge \hat y_f\not\sqsubseteq d], \label{eq:ungr}\\
\mathrm{Inc}(\hat y) &= \mathbf{1}\big[\,|\mathrm{CVSS}(\hat y_{\text{vec}})-\hat y_{\text{score}}|>0.05\,\big]. \label{eq:inc}
\end{align}
$\mathrm{Inv}$ flags non-existent CWEs, malformed vectors and out-of-range scores. $\mathrm{Ungr}$ flags a vendor or product name that does not occur in the record text $d$. $\mathrm{Inc}$ flags a score that contradicts the model's own vector.
Here $\sqsubseteq$ denotes substring containment after slug normalisation, and $\mathrm{CVSS}$ is the exact v3.1 base equation with its round-up rule \citep{first2019cvss}.

\subsection{Aggregation Operators}
Let $\hat y_1,\dots,\hat y_n$ be member answers for one field of one record. Abstentions are non-votes, but they still count in the denominator, so silence is never consent:
\begin{equation}\label{eq:count}
\kappa(v)=\textstyle\sum_{m=1}^{n}\mathbf{1}[\hat y_m=v],\qquad v^\dagger=\arg\max_{v\neq\bot}\kappa(v).
\end{equation}
If the arg max is not unique, every operator returns $\bot$. Otherwise:
\begin{align}
\mathsf{Plur}&=v^\dagger, \qquad
\mathsf{Maj}=v^\dagger\ \text{if}\ \kappa(v^\dagger)>n/2\ \text{else}\ \bot,\label{eq:maj}\\
\mathsf{kOfN}_k&=v^\dagger\ \text{if}\ \kappa(v^\dagger)\ge k\ \text{else}\ \bot, \qquad
\mathsf{Unan}=\mathsf{kOfN}_n.\label{eq:kofn}
\end{align}
\emph{Verified majority} applies a filter $\Phi$ before voting, replacing with $\bot$ every member value that fails Eqs.~\eqref{eq:inv}--\eqref{eq:inc} (for $\mathrm{Inc}$, both the member's score and its vector are removed), and then applies $\mathsf{Maj}$ with the original $n$: $\mathsf{VMaj}=\mathsf{Maj}\circ\Phi$. \emph{Metric-wise voting} (T2 only) takes a plurality over each of the eight CVSS metrics among members with a valid vector, requires more than $n/2$ valid vectors, and recomputes the score as $\mathrm{CVSS}(\cdot)$ of the voted vector. \emph{Judge fusion} passes the record and the $n$ anonymised candidate answers to an LLM, which is instructed to keep a field only if candidates agree and never to introduce an unproposed value. We record the rate at which it does so anyway (the \emph{novel-value} rate).

\paragraph{Compute-matched self-consistency.} For a single model $M$ we draw $S=n$ samples at temperature $\tau$ and apply the same operators. An ensemble and its self-consistency counterpart therefore make the same number of calls. Their difference isolates the contribution of using \emph{different} models \citep{wang2023selfconsistency}.

\subsection{When Can Voting Remove a Fabrication?}

\begin{proposition}[Majority suppression of a specific fabricated value]\label{prop:hoeffding}
Fix a record and field with gold value $y^*$, and assume member outputs are independent given the record. For a wrong value $v\neq y^*$ let $\pi_v=\frac1n\sum_m\Pr[\hat y_m=v]$. If $\pi_v<\tfrac12$,
\begin{equation}
\Pr[\mathsf{Maj}=v]\ \le\ \exp\!\big(-2n(\tfrac12-\pi_v)^2\big),\qquad
\Pr[\mathsf{Maj}\ \text{fabricates}]\ \le\ \textstyle\sum_{v\neq y^*}\exp\!\big(-2n(\tfrac12-\pi_v)^2\big).
\end{equation}
\end{proposition}
\begin{proof}
$\kappa(v)=\sum_m \mathbf{1}[\hat y_m=v]$ is a sum of $n$ independent Bernoulli variables with mean $n\pi_v$. $\mathsf{Maj}=v$ requires $\kappa(v)>n/2$, i.e.\ $\kappa(v)-n\pi_v>n(\tfrac12-\pi_v)$. Hoeffding's inequality \citep{hoeffding1963} gives the first bound, and a union bound over $v$ gives the second.
\end{proof}

\begin{remark}[Total versus identical fabrication]\label{rem:identical}
The bound depends on $\pi_v$, the probability mass that members place on \emph{one specific} wrong value, and not on the total fabrication rate $\sum_v\pi_v$. Members that fabricate often but \emph{differently} are filtered efficiently. Members that converge on the same ``popular'' wrong value (e.g.\ score 9.8 or 7.5, CWE-20 or CWE-79) are not. The independence assumption fails in exactly this case \citep{kim2025correlated}. We therefore measure, for every model pair, the \emph{identical-fabrication rate}: the share of records on which both models fabricate \emph{the same} value, also reported conditional on both fabricating. This is the quantity that predicts whether an ensemble can help.
\end{remark}

\begin{proposition}[Verification can only trade coverage for safety]\label{prop:verified}
For every record and field, if $\mathsf{VMaj}\neq\bot$ then $\mathsf{VMaj}=\mathsf{Maj}$. Consequently $\Fab(\mathsf{VMaj})\le\Fab(\mathsf{Maj})$, $\Acc(\mathsf{VMaj})\le\Acc(\mathsf{Maj})$ and $\Cov(\mathsf{VMaj})\le\Cov(\mathsf{Maj})$. Likewise, $\mathsf{kOfN}_{k'}\neq\bot\Rightarrow\mathsf{kOfN}_{k'}=\mathsf{kOfN}_{k}$ for $k'\ge k$, so the agreement threshold traces a nested risk--coverage curve.
\end{proposition}
\begin{proof}
$\Phi$ only replaces values with $\bot$, so $\kappa_\Phi(v)\le\kappa(v)$ for every $v$, while the threshold $n/2$ is unchanged. If $\mathsf{VMaj}=v$ then $\kappa(v)\ge\kappa_\Phi(v)>n/2$. Since at most one value can exceed $n/2$, $v$ is the unique arg max before filtering, so $\mathsf{Maj}=v$. The rate inequalities follow because the answered set of $\mathsf{VMaj}$ is a subset of that of $\mathsf{Maj}$, with identical answers. The $\mathsf{kOfN}$ case is identical, with threshold $k'$.
\end{proof}

Proposition~\ref{prop:verified} turns H4 into a guarantee about the \emph{direction} of the effect. The empirical question is only how much coverage verification costs.

\subsection{Diversity Statistics}
For members $a,b$ with correctness indicators on the same records, let $N^{11},N^{00},N^{10},N^{01}$ count both-correct, both-wrong and the two discordant cells. We report the disagreement measure $(N^{10}+N^{01})/N$ and Yule's $Q=(N^{11}N^{00}-N^{01}N^{10})/(N^{11}N^{00}+N^{01}N^{10})$ \citep{kuncheva2003}, together with the identical-fabrication rate of Remark~\ref{rem:identical} and the oracle (``any member correct'') upper bound.

\subsection{Statistical Testing}
Every system is compared with the \emph{best single model}, selected on the development split by accuracy (ties broken by lower fabrication) and evaluated on the disjoint test split. We use McNemar's exact test \citep{mcnemar1947} on paired correctness and, separately, on paired non-fabrication. Let $b$ be the number of records where the reference fails and the system succeeds, and $c$ the converse; the two-sided $p$-value is $2\,\mathrm{Binom}(\min(b,c);\,b{+}c,\tfrac12)$, capped at 1. $p$-values are Holm-corrected within each (task, field) family \citep{holm1979}. Percentile bootstrap 95\% confidence intervals use 1000 resamples of test records \citep{efron1993}.

% ============================================================================
\section{Methodology}\label{sec:method}

\subsection{Dataset Collection and Profile}
Records are drawn from the NVD CVE API~2.0 \citep{nvd2024api} and kept only if they carry a CVSS v3.1 base metric (NVD's own primary assessment where available). Two strata are sampled:
\begin{itemize}[leftmargin=1.5em]
\item \textbf{Recent (post-cutoff):} $N_r=1000$ records published between \TBD{start} and \TBD{end}, a window beginning after the latest documented training cutoff of every model in the registry. For T1, the correct behaviour on this stratum is to abstain, so any answer that matches the gold value is either memorised leakage or a lucky guess.
\item \textbf{Control (pre-cutoff):} $N_c=300$ records published in 2021, which every model may have seen in training. This mirrors the temporal control of \citet{anon2026grounding}.
\end{itemize}
Each record is stored as $R_i=(c_i,d_i,s^*_i,\mathbf{v}^*_i,\mathrm{cwe}^*_i,\mathrm{cpe}^*_i)$. A deterministic hash of the identifier assigns 20\% of records to a development split, which is used only to select the best single model, and 80\% to the test split, on which all results are reported. Table~\ref{tab:dataset} profiles both strata, including the share of records with a gradable CWE and with CPE configurations. The latter limits how many records can be graded for vendor and product.

\resultstable{tab_dataset}{Corpus profile per stratum: size, CVSS v3.1 score mean / median / SD, NIST severity bands, and the share of records with a gradable CWE and with CPE configurations.}{tab:dataset}

\subsection{Tasks and Prompts}
\textbf{T1 (parametric recall)} gives only the identifier and asks for the NVD CVSS v3.1 base score and primary CWE. \textbf{T2 (record-grounded inference)} gives the identifier and the description, with all labels removed, and asks for the base vector, base score, primary CWE, vendor and product. It also states that vendor and product must be taken from the description. Both prompts request a single JSON object and explicitly permit \texttt{null} (Appendix~\ref{app:prompts}). Permitting abstention is a deliberate change from the score-only prompt of \citet{anon2026grounding}, and \TBD{an ablation with \texttt{abstain\_instruction=false}} quantifies its effect.

\subsection{Model Registry, Execution Parameters, and Hardware}
The ensemble deliberately mixes vendors, training corpora and serving modes, since Remark~\ref{rem:identical} predicts that shared provenance increases identical fabrications \citep{kim2025correlated}. It retains the open-weight models of the earlier study for continuity.

\begin{table}[htbp]\centering\small
\caption{Model registry. Exact model versions and access dates are recorded in the released run configuration.}\label{tab:models}
\begin{adjustbox}{max width=\textwidth}\begin{tabular}{@{}llll@{}}\toprule
Member & Family / provider & Serving & Max output tokens \\ \midrule
Mistral-7B-Instruct-v0.2 \citep{jiang2023mistral} & Mistral (open) & vLLM, local GPU & 200 \\
Qwen2.5-7B-Instruct \citep{qwen2024qwen25} & Alibaba (open) & vLLM, local GPU & 200 \\
Llama-3.1-8B-Instruct \citep{grattafiori2024llama3} & Meta (open) & vLLM, local GPU & 200 \\
GPT-5-mini & OpenAI (hosted, reasoning) & API & 4000$^\dagger$ \\
\TBD{Claude model} & Anthropic (hosted) & API & 300 \\
\TBD{Gemini model} & Google (hosted) & API & 2000 \\
\bottomrule
\multicolumn{4}{@{}l}{\footnotesize $^\dagger$Reasoning tokens count against the completion budget; see text.}
\end{tabular}\end{adjustbox}
\end{table}

Open-weight models are served with vLLM \citep{kwon2023vllm} on \TBD{GPU}. Ensemble members answer once with greedy decoding ($T{=}0$), so that the comparison isolates model diversity from sampling noise. Self-consistency draws $S{=}n$ samples at $\tau{=}0.7$ from \TBD{Qwen2.5-7B and GPT-5-mini}. Hosted reasoning models only accept their default temperature and are run with reasoning effort \texttt{low}. Their completion budget is set to 4000 tokens because hidden reasoning tokens count against it: under a short budget the visible reply can be empty with \texttt{finish\_reason=length}, which a parser would read as a refusal. The harness logs the finish reason and empty-reply rate of every call (Table~\ref{tab:diag}). This lets us test one candidate explanation for the universal GPT-5-mini abstention reported by \citet{anon2026grounding}.

All calls are issued concurrently with per-backend concurrency limits and exponential back-off. Responses are appended to per-model JSONL logs keyed by (model, task, CVE, sample, temperature, prompt version), so interrupted runs resume without re-querying and every number in this paper can be regenerated from the logs.

\subsection{Evaluation Pipeline Algorithms}

\begin{algorithm}[htbp]
\caption{Parallel fan-out (phase 1)}\label{alg:fanout}
\begin{algorithmic}[1]
\Require records $\mathcal{D}$, tasks $\mathcal{T}$, members $\mathcal{M}$, self-consistency set $\mathcal{M}_{sc}\subseteq\mathcal{M}$, $S$, $\tau$
\ForAll{$T\in\mathcal{T}$, $R_i\in\mathcal{D}$ \textbf{in parallel}}
  \State $p\gets\textsc{Prompt}_T(R_i)$ \Comment{labels stripped; null permitted}
  \ForAll{$M_m\in\mathcal{M}$ \textbf{in parallel}}
     \State $r_{m,i}\gets M_m(p;\,T{=}0)$ \textbf{unless cached}; log text, finish reason, tokens, latency
  \EndFor
  \ForAll{$M\in\mathcal{M}_{sc}$, $s=1..S$ \textbf{in parallel}}
     \State $r^{(s)}_{M,i}\gets M(p;\,\tau)$ \textbf{unless cached}
  \EndFor
\EndFor
\end{algorithmic}
\end{algorithm}

\begin{algorithm}[htbp]
\caption{Verified consensus for one record (phase 3)}\label{alg:vmaj}
\begin{algorithmic}[1]
\Require raw replies $r_1..r_n$, fields $\mathcal{F}_T$, record text $d$, CWE catalogue $\mathcal{C}$
\For{$m=1..n$}
  \State $\hat y_m\gets\textsc{Normalise}(\textsc{ParseJSON}(r_m),\mathcal{C})$ \Comment{invalid values kept in $\mathcal{X}_f$}
  \ForAll{$f\in\mathcal{F}_T$}
    \If{$\mathrm{Inv}(\hat y_{m,f})\ \vee\ \mathrm{Ungr}(\hat y_{m,f},d)$} $\hat y_{m,f}\gets\bot$ \EndIf
  \EndFor
  \If{$\mathrm{Inc}(\hat y_m)$} $\hat y_{m,\text{score}}\gets\bot;\ \hat y_{m,\text{vec}}\gets\bot$ \EndIf
\EndFor
\ForAll{$f\in\mathcal{F}_T$}
  \State $v^\dagger\gets$ unique $\arg\max_{v\neq\bot}\kappa_f(v)$, else $\bot$
  \State $\hat y_f\gets v^\dagger$ \textbf{if} $\kappa_f(v^\dagger)>n/2$ \textbf{else} $\bot$ \Comment{abstentions count in $n$}
\EndFor
\State \Return $\hat y$ \Comment{$\bot$ fields route to analyst review}
\end{algorithmic}
\end{algorithm}

\begin{algorithm}[htbp]
\caption{LLM-judge fusion (phase 2)}\label{alg:judge}
\begin{algorithmic}[1]
\Require record $R_i$, member answers $\hat y_1..\hat y_n$, judge model $J$
\State $q\gets\textsc{JudgePrompt}(R_i,\ \text{anonymised}\ \hat y_1..\hat y_n)$ \Comment{``never introduce an unproposed value''}
\State $\hat y^J\gets\textsc{Normalise}(\textsc{ParseJSON}(J(q;\,T{=}0)))$
\State $\mathrm{novel}\gets\exists f:\ \hat y^J_f\neq\bot\wedge \hat y^J_f\notin\{\hat y_{m,f}\}_{m}$ \Comment{fusion-introduced fabrication}
\State \Return $\hat y^J,\ \mathrm{novel}$
\end{algorithmic}
\end{algorithm}

% ============================================================================
\section{Experimental Results}\label{sec:results}
All results are on the test split. Systems are grouped as single models, compute-matched self-consistency (SC) and parallel ensembles (Ens). All rates are percentages of gradable records.

\subsection{Response Diagnostics}
Before comparing systems, we check that abstentions are genuine model decisions rather than artefacts. Table~\ref{tab:diag} reports, per model and task, the share of empty replies, replies truncated by the token budget, and replies that parsed as JSON.

\resultstable{tab_diagnostics}{Response diagnostics per model and task: empty replies, truncation (\texttt{finish\_reason=length}), JSON parse success and mean output tokens.}{tab:diag}

\begin{keyresult}
\TBD{Do empty or truncated replies explain any model's abstention rate? In particular, report GPT-5-mini's empty-reply rate under the 4000-token budget and compare its abstention with the 100\% reported earlier.}
\end{keyresult}

\subsection{RQ1: Parametric Recall}
Table~\ref{tab:recall} and Figure~\ref{fig:failure-recall} compare single models, self-consistency and ensembles when only the identifier is given. Table~\ref{tab:strata} splits the result by stratum.

\resultstable{tab_recall_main}{T1 (identifier only): accuracy, fabrication and abstention (\%) for CVSS score and CWE, with macro selective accuracy and invalid-value rate.}{tab:recall}
\resultsfig{fig_failure_recall}{T1 failure-mode decomposition for the CVSS score: correct, abstained and fabricated share per system.}{fig:failure-recall}
\resultstable{tab_strata}{T1 by stratum. On the post-cutoff stratum, abstention is the only safe answer.}{tab:strata}

\begin{keyresult}
\TBD{RQ1 answer: post-cutoff fabrication of the best single model vs.\ Ens majority / verified majority (with CIs); whether plurality inherits shared fabrications (H2); control-stratum accuracy.}
\end{keyresult}

\subsection{RQ2: Record-Grounded Inference}
Table~\ref{tab:record} reports macro-averaged outcomes and the ground-truth-free flags for T2. Tables~\ref{tab:fields-acc}--\ref{tab:fields-fab} give per-field accuracy and fabrication.

\resultstable{tab_record_main}{T2 (identifier + description), macro over fields: accuracy, fabrication, abstention, selective accuracy, and the share of invalid, ungrounded (vendor/product) and self-inconsistent (score) outputs.}{tab:record}
\resultstable{tab_record_fields_acc}{T2 per-field accuracy (\%).}{tab:fields-acc}
\resultstable{tab_record_fields_fab}{T2 per-field fabrication (\%).}{tab:fields-fab}
\resultsfig{fig_failure_record}{T2 failure-mode decomposition (macro over fields).}{fig:failure-record}

\begin{keyresult}
\TBD{RQ2 answer: fields where the ensemble helps most and least; reduction in ungrounded vendor/product and inconsistent vector/score; whether metric-wise voting beats exact-vector voting.}
\end{keyresult}

\subsection{RQ3: Diversity versus Compute}
Table~\ref{tab:diversity} summarises pairwise error diversity, and Figure~\ref{fig:samefab} shows identical-fabrication rates for every model pair. The SC rows of Tables~\ref{tab:recall} and~\ref{tab:record} give the compute-matched comparison.

\resultstable{tab_diversity}{Mean pairwise diversity among ensemble members: disagreement, Yule's $Q$, both-fabricate rate, identical-fabrication rate (unconditional and given both fabricate), and the oracle upper bound.}{tab:diversity}
\resultsfig[0.48]{fig_samefab_recall_cvss_score}{Identical fabrications on T1 score: share of joint fabrications in which two models output the same wrong value.}{fig:samefab}

\begin{keyresult}
\TBD{RQ3 answer: ensemble vs.\ SC at equal call budget; identical-fabrication rate for cross-vendor vs.\ same-model pairs (H3); whether Remark~\ref{rem:identical} predicts which fields benefit.}
\end{keyresult}

\subsection{RQ4: Aggregation Rule, Ensemble Size and Operating Point}
Figure~\ref{fig:size} plots accuracy and fabrication against ensemble size, averaged over all subsets of each size. Figure~\ref{fig:riskcov} plots the risk--coverage curve traced by the agreement threshold $k$, with and without pre-vote verification, against each single model's operating point. With strict majority, even ensemble sizes require a larger supermajority than the neighbouring odd sizes, which produces a characteristic odd--even oscillation.

\resultsfig{fig_size_recall_cvss_score}{Ensemble size $k$ vs.\ accuracy and fabrication (T1 score). Lines: mean over all $\binom{n}{k}$ subsets; bands: min--max.}{fig:size}
\resultsfig{fig_size_record_cwe}{Ensemble size $k$ vs.\ accuracy and fabrication (T2 CWE).}{fig:size-cwe}
\resultsfig[0.5]{fig_riskcov_record_cwe}{Risk--coverage for T2 CWE: agreement threshold $k$ (labelled) with and without verification; markers are single models.}{fig:riskcov}

\begin{keyresult}
\TBD{RQ4 answer: recommended rule and $k$; coverage cost of verification (Proposition~\ref{prop:verified} guarantees no fabrication increase); judge novel-value rate (H5).}
\end{keyresult}

\subsection{Statistical Significance}
\resultstable{tab_mcnemar}{McNemar exact tests against the best single model (chosen on the development split). $\Delta$ in percentage points; $b/c$ = discordant counts; Holm-corrected $p$ within each task--field family. Full table in the released results.}{tab:mcnemar}

\subsection{Cost and Latency}
\resultstable{tab_cost}{Mean tokens per CVE, mean wall-clock latency when members run in parallel vs.\ sequentially, and API cost per 1000 CVEs (hosted models; local inference reported as \$0).}{tab:cost}

\begin{keyresult}
\TBD{cost multiplier of the recommended ensemble over the best single model; parallel latency overhead.}
\end{keyresult}

% ============================================================================
\section{Discussion}\label{sec:discussion}

\subsection{Consensus Filters Idiosyncratic, Not Shared, Fabrications}
Proposition~\ref{prop:hoeffding} predicts that the benefit of voting is governed by how members distribute their errors, not by how often they err. A heterogeneous ensemble whose members fabricate frequently but idiosyncratically can be safer than any member. An ensemble of individually stronger but correlated models may not be. \TBD{Relate the measured identical-fabrication rates (Table~\ref{tab:diversity}) to the per-field ensemble gains (Tables~\ref{tab:fields-acc}--\ref{tab:fields-fab}); state whether fields with high identical-fabrication rates, e.g.\ generic CWEs, are where voting fails.} This is consistent with the marginal ensemble gains reported for CVSS-from-description \citep{jafarikhah2025description} and with the cross-model error correlation documented by \citet{kim2025correlated}.

\subsection{Abstention Is the Price, and Usually Worth Paying}
Strict majority and verification convert disagreements into abstentions. In a vulnerability pipeline an abstention is an auditable routing event: the record goes to an analyst or waits for NVD enrichment. A fabrication silently alters prioritisation. Proposition~\ref{prop:verified} guarantees that verification cannot make fabrication worse, so the operational choice reduces to selecting a point on the risk--coverage curve (Figure~\ref{fig:riskcov}) that matches analyst capacity. \TBD{Quantify: the coverage at which the ensemble matches or beats the best single model's selective risk.}

\subsection{Diversity Versus Compute}
Self-consistency spends the same number of calls on one model. Its samples share training data, tokenizer and alignment, so its fabrications are expected to be more identical than those of a cross-vendor ensemble. \TBD{Report whether the SC rows match the ensemble rows; if they do, the cheaper single-provider deployment is preferable.}

\subsection{Relation to Retrieval Grounding}
Ensembles do not replace grounding. When an authoritative record exists, retrieval plus verification already achieves zero fabrication \citep{anon2026grounding}, and voting over grounded answers adds cost without benefit. The ensemble's role is the gap that grounding cannot fill: inferring metadata for records not yet enriched, and screening descriptions before enrichment. The ground-truth-free flags (Eqs.~\eqref{eq:inv}--\eqref{eq:inc}) carry the dual-gate philosophy of the earlier study into this ungrounded setting.

\subsection{Operational Guidance}
\TBD{Final recommendations, e.g.: (i) use an odd number of cross-vendor members; (ii) apply deterministic verification before voting; (iii) prefer strict majority over plurality for recall-type queries; (iv) use judge fusion only with novel-value checking; (v) route $\bot$ to human review.}

% ============================================================================
\section{Limitations}\label{sec:limitations}
\begin{enumerate}[leftmargin=1.6em]
\item \textbf{Ground-truth quality.} NVD labels are themselves analyst judgements, and CVSS vectors from different assessors can legitimately differ. Exact-vector accuracy is therefore a strict criterion, and some ``fabrications'' may be defensible alternative assessments.
\item \textbf{Partial gradability.} Vendor and product can only be graded where NVD has published CPE configurations, and CWE only where a specific CWE is assigned. Graded subsets may differ systematically from the full corpus (Table~\ref{tab:dataset}).
\item \textbf{Training cutoffs.} The post-cutoff window relies on documented model cutoffs. Undisclosed continued training or retrieval inside hosted models could leak post-cutoff records.
\item \textbf{Model and prompt scope.} Results hold for the registry, prompts and decoding settings evaluated. Model versions behind hosted endpoints can change over time. All prompts are in English and use a single JSON schema.
\item \textbf{Independence in Proposition~\ref{prop:hoeffding}.} The bound assumes conditionally independent members. Measured correlation is reported, but the bound is not corrected for it.
\item \textbf{Judge contamination.} The judge model is also an ensemble member, which may bias fusion towards its own answer.
\item \textbf{Adversarial inputs.} We do not evaluate descriptions crafted to induce coordinated fabrication across models, such as prompt injection embedded in advisories.
\end{enumerate}

\section{Future Work}\label{sec:future}
\begin{enumerate}[leftmargin=1.6em]
\item Correlation-aware aggregation that down-weights members whose fabrications historically coincide.
\item Cascades that query a second model only when the first is uncertain, reducing the cost of consensus \citep{chen2023frugalgpt}.
\item Combining ensembles with retrieval: voting over grounded answers when the record exists, and voting over inferences otherwise.
\item CVSS v4.0 and CNA-assigned (non-NVD) labels; multilingual advisories.
\item Adversarial robustness of consensus against crafted advisories.
\end{enumerate}

\section{Conclusion}\label{sec:conclusion}
We asked whether parallel multi-LLM consensus prevents fabrication in CVE record analysis. The answer depends on whether models fabricate the same values, so we measured that directly, alongside accuracy, fabrication and abstention for every single model, every ensemble rule and a compute-matched self-consistency baseline. \TBD{Two- to three-sentence summary of findings for RQ1--RQ4.} The harness and all per-query outputs are released so that the same comparison can be repeated as models change.

% ============================================================================
\section*{Declarations}
\paragraph{Funding.} \TBD{}
\paragraph{Competing interests.} The authors declare no competing interests.
\paragraph{Ethics approval and consent to participate.} Not applicable. The study uses only public, non-personal vulnerability metadata from the NVD and involves no human participants.
\paragraph{Consent for publication.} Not applicable.
\paragraph{Data availability.} The corpus is drawn from the public NVD CVE API~2.0. The sampled identifiers, labels and every raw model response underlying the tables and figures will be deposited in a public repository \TBD{URL}.
\paragraph{Code availability.} The \texttt{cvebench} harness (data collection, model backends, parallel runner, aggregation, analysis and figure generation) will be released \TBD{URL}. API credentials are read from environment variables and never stored.
\paragraph{Author contributions.} \TBD{} (withheld in the blinded manuscript).

\bibliographystyle{plainnat}
\bibliography{references}

% ============================================================================
\appendix
\section{Prompts}\label{app:prompts}
\small
\noindent System message (all tasks): \emph{You are a cybersecurity vulnerability analyst. Reply with a single JSON object and nothing else.}

\medskip\noindent\textbf{T1 (recall).}
\begin{verbatim}
Give the NVD CVSS v3.1 base score and the primary CWE for {cve_id}.
If you cannot determine a field with confidence, set it to null instead of guessing.
Return exactly this JSON:
{"cvss_score": <number 0.0-10.0 or null>, "cwe": "<CWE-ID like CWE-79 or null>"}
\end{verbatim}

\noindent\textbf{T2 (record).}
\begin{verbatim}
Analyse the CVE record below.
CVE ID: {cve_id}
Description: {description}
Infer the CVSS v3.1 base vector and base score, the primary CWE, and the
affected vendor and product. Vendor and product must be taken from the
description; do not invent them.
If you cannot determine a field with confidence, set it to null instead of guessing.
Return exactly this JSON:
{"cvss_vector": "<CVSS:3.1/AV:_/AC:_/PR:_/UI:_/S:_/C:_/I:_/A:_ or null>",
 "cvss_score": <number or null>, "cwe": "<CWE-ID or null>",
 "vendor": "<string or null>", "product": "<string or null>"}
\end{verbatim}

\noindent\textbf{Judge.}
\begin{verbatim}
Several independent analysts answered the same question about a CVE. Their
answers may contain fabricated values.
{record or question}
Candidate answers:
Analyst A: {...}
...
Produce one final answer. Keep a field only if it is supported by the record
(when given) and by the candidates' agreement; otherwise set it to null. Never
introduce a value that no candidate proposed.
Return exactly this JSON with the same keys as the candidates: {...}
\end{verbatim}

\section{Reproducibility}\label{app:repro}
\begin{verbatim}
python -m cvebench.data.fetch_nvd --start <recent_start> --end <recent_end> \
       --n 1000 --stratum recent --out data/cves.jsonl
python -m cvebench.data.fetch_nvd --start 2021-01-01 --end 2021-12-31 \
       --n 300 --stratum control --out data/cves.jsonl --append
python -m cvebench.data.fetch_cwe --out data/cwe_ids.txt
python -m cvebench.cli all --config configs/main.yaml
cd paper && latexmk -pdf main.tex
\end{verbatim}
The CVSS v3.1 calculator is unit-tested against reference vectors from the FIRST specification, and the full pipeline is exercised end-to-end on synthetic records with simulated models before any real run.
\end{document}
